\chapter{Related Work}\label{ch:related}

This chapter reviews existing literature related to streaming communication in microservice-based systems, with a particular focus on topics that are central to this thesis: streaming communication technologies and interaction models, architectural integration in adaptive microservice middleware, policy enforcement under continuous data exchange, resilience mechanisms for streaming workflows, and validation practices with respect to scalability and performance.

\section{Streaming communication technologies and interaction models}

Policy-aware data exchange systems constrain communication choices because the data path has to remain compatible with orchestration and compliance boundaries. DYNAMOS, the case-study system introduced in \Cref{ch:background}, is deployed on Kubernetes, uses gRPC for inter-service communication, and uses RabbitMQ for asynchronous messaging~\cite{stutterheim2023thesis,stutterheim2023icsa}. This constrains the case-study design space to technologies that are compatible with containerised microservices and cloud-native deployment models. The primary question addressed in this body of work is which streaming communication mechanisms are suitable for large-scale or continuous data exchange, and under which conditions direct streaming RPC should be preferred over broker-based approaches.

gRPC provides native support for streaming through client-streaming, server-streaming, and bidirectional streaming RPCs. These interaction patterns differ in message flow control, lifetime of connections, buffering behaviour, and error propagation semantics. In the context of DYNAMOS, the choice of streaming pattern is not merely an API concern but has architectural consequences, as it affects memory usage, cancellation behaviour, and the ability to overlap communication and computation. Stutterheim~\cite{stutterheim2023thesis} identifies client-streaming as a promising mechanism for transferring large datasets incrementally, reducing peak memory pressure at the cost of increased handling complexity.

However, the use of long-lived streaming connections introduces operational challenges that are largely absent in unary request-response interactions. Starck and Ghofrani~\cite{starck2023streaminglb} show that gRPC streaming complicates load balancing in Kubernetes environments, as persistent connections can cause traffic to become unevenly distributed across replicas. This observation is particularly relevant for systems such as DYNAMOS that rely on horizontal scaling and ephemeral execution units. Complementary benchmarking work demonstrates that the performance characteristics of streaming gRPC are sensitive to runtime configuration and virtualisation overhead, with measurable effects on latency and resource utilisation in containerised environments \cite{sodankoor2026grpcbench}.

An alternative to direct service-to-service streaming is to externalise the data plane to a messaging or event streaming platform. Broker-based communication decouples producers and consumers in time, delegates buffering and persistence to the broker, and can simplify fan-out scenarios. Comparative benchmarks of Kafka, RabbitMQ, NATS, and ActiveMQ Artemis show that different brokers dominate under different workload characteristics, reinforcing that broker choice is highly workload dependent \cite{ibrahim2026brokers}. A cloud-native literature review comparing Kafka, Pulsar, and RabbitMQ further highlights trade-offs between throughput, reliability, operational complexity, and Kubernetes integration \cite{nuikka2021cloudnative}. While these studies are not conducted in the context of policy-driven middleware, they provide structured criteria that are directly applicable when evaluating broker-based alternatives for DYNAMOS.

Several studies also emphasise that streaming design is influenced not only by technology choice but by the underlying communication model. DataXc contrasts push-based streaming approaches with pull-based designs and reports improved latency stability and throughput under consumer-controlled flow \cite{coviello2022dataxc}. Although the application domain differs, the distinction between producer-driven and consumer-driven streaming is directly relevant to backpressure management and resource utilisation in DYNAMOS-compatible workflows.

Overall, existing literature provides a broad comparison of streaming technologies and interaction patterns, but these comparisons are typically conducted in generic microservice or stream-processing settings. Existing studies are still useful methodologically because they emphasise identical workloads, common measurement points, latency distribution, throughput, resource pressure, and recovery behaviour as recurring benchmark dimensions \cite{karimov2018benchmarking,Henning2021ScalabilityMetrics,systematicstreaming,ciftci2025ipcmechanisms}. The gap addressed in this thesis is therefore not the invention of a new benchmarking methodology, but the application of those dimensions to dynamically composed, policy-driven middleware with ephemeral execution units.

\section{Architectural integration of streaming in adaptive microservice middleware}

Beyond the selection of a streaming mechanism, integrating streaming communication into an adaptive microservice architecture raises architectural questions related to modularity, orchestration, and lifecycle management. DYNAMOS is explicitly designed around dynamic microservice composition, where execution chains are generated and deployed per request based on policy evaluation and system state \cite{stutterheim2023thesis,stutterheim2023icsa}. These chains are short-lived by design, which improves isolation and compliance guarantees but complicates the management of long-lived communication channels.

Communication in DYNAMOS is abstracted through a sidecar-based architecture, decoupling communication concerns from core microservice logic. Figueira and Coutinho~\cite{figueira2024selfadaptive} show that sidecars can support runtime adaptation of communication behaviour in self-adaptive Kubernetes architectures, and discuss sidecar runtimes such as Dapr that unify service invocation, pub/sub, and state management behind a protocol-agnostic interface. This line of work establishes that heterogeneous communication paradigms, including streaming and messaging, can be hidden behind a stable sidecar interface. What it does not address is dynamic chain composition driven by policy evaluation: sidecar runtimes assume relatively static service topologies, so the combination of streaming transfer with ephemeral, policy-generated execution chains remains open. That combination is precisely the setting of this thesis.

\section{Policy enforcement and compliance under continuous data exchange}

Policy enforcement is a defining characteristic of DYNAMOS, where access rights, execution locations, and data-flow constraints are evaluated prior to execution and enforced through dynamically composed microservice chains \cite{stutterheim2023thesis,stutterheim2023icsa}. In the current architecture, policy enforcement is implicitly scoped to bounded request-response interactions. Introducing streaming communication challenges this assumption, as data transfer may span longer periods during which policies, system state, or execution context may change.

The access-control literature has a name for this distinction. Park and Sandhu~\cite{park2004ucon} introduce usage control (UCON) as a generalisation of access control in which authorisation is not a one-time decision at request time: decisions can be re-evaluated during use, attributes can change while access is ongoing, and access can be revoked mid-use. Hariri et al.~\cite{hariri2023uconplus} extend this model with an architecture for continuous authorisation in which policy re-evaluation is triggered by attribute changes during an active session. In the data-exchange domain, dataspace connectors such as the IDS Dataspace Connector operationalise usage-control policies for inter-organisational transfer~\cite{idsa2023usagecontrol}. This vocabulary matters for streaming middleware because it separates two enforcement models: pre-decision enforcement, where a request is approved once before execution, and ongoing enforcement, where an active transfer remains subject to re-evaluation.

Neumaier et al.~\cite{neumaier2021policyaware} make the same distinction concrete for decentralised dataset exchange, proposing continuous monitoring and runtime validation against machine-readable policies. While their work does not explicitly target streaming RPC, the distinction between one-time policy clearance and ongoing compliance monitoring becomes critical when data is exchanged incrementally over long-lived connections.

At the communication layer, several studies propose enforcing security and access control through proxies or sidecars. Thiyagarajan et al.~\cite{thiyagarajan2025grpcsecurity} present a proxy-based approach for securing gRPC communication using mTLS, JWT authentication, and RBAC. Although the focus is on security rather than data usage policy, the architectural approach is relevant, as it suggests that enforcement can be applied at communication boundaries and potentially at message granularity.

Despite these contributions, there is limited guidance on how usage-control-style ongoing enforcement should interact with long-lived streaming connections in systems that dynamically compose execution chains. In particular, existing work does not address how policy violations should be detected or handled mid-stream, or how streams should be terminated or adapted when compliance conditions change. These limitations highlight open challenges when introducing streaming into policy-driven middleware such as DYNAMOS, and they motivate the explicit enforcement-model boundary drawn in \Cref{ch:dynamos-integration}.

\section{Resilience, backpressure, and fault handling in streaming workflows}

Streaming-based communication introduces operational challenges that are less pronounced in request-response systems. Long-lived data flows must tolerate partial failures, adapt to fluctuating load, and prevent uncontrolled resource growth through effective backpressure mechanisms. These challenges are compounded in DYNAMOS by dynamic microservice composition and ephemeral execution units.

The most directly relevant prior work is fault-recovery benchmarking. Vogel et al.~\cite{vogel2024faultRecovery} benchmark fault recovery in Flink, Kafka Streams, and Spark Structured Streaming by injecting worker failures in a cloud-native testbed and measuring recovery behaviour over a fixed observation window. They argue that recovery has been under-measured relative to clean throughput, and that meaningful recovery evaluation requires several complementary measures: the time until output resumes, the time until throughput is restored to its pre-failure level, the backlog accumulated during the outage, and the correctness cost of replay under at-least-once delivery. This thesis adopts that structure at the transport level rather than the framework level: the recovery scenario in \Cref{ch:transport-benchmark} injects a component restart and reports resumption time, throughput restoration, accumulated backlog, and post-failure duplicates and ordering violations per transport mapping.

Beyond recovery benchmarking, reactive microservice architectures emphasise asynchronous communication, non-blocking execution, and fault isolation as resilience mechanisms under variable load \cite{rasheedh2021reactive}, and Guntupalli and Vamshi~\cite{guntupalli2023highvolume} show that event-driven architectures with messaging systems can mitigate bottlenecks under sustained data rates. From a systems perspective, Gan and Delimitrou~\cite{gan2018architectural} demonstrate that communication overhead and tail latency dominate the performance of large-scale microservice deployments, so resilience and scalability issues often emerge from inter-service communication rather than computation.

Existing literature therefore provides both principles (fault isolation, explicit flow control) and measurement practice (fault-injection benchmarking with multi-dimensional recovery metrics), but applies them either to request-oriented architectures or to stream-processing frameworks with static topologies. There is limited work examining how these mechanisms behave for transport-level streaming in dynamically composed microservice chains.

\section{Validation practices, scalability, and performance evaluation}

Empirical evaluation of streaming communication often focuses on comparing streaming and non-streaming interaction styles in terms of throughput, latency, and resource usage. Studies focused specifically on gRPC streaming show that streaming RPCs outperform unary calls for large payloads and continuous data flows by enabling incremental processing and reducing buffering overhead \cite{roohitavaf2019logplayer,sodankoor2026grpcbench}. Complementary protocol-level comparisons between REST and gRPC further establish that binary, multiplexed transports outperform request-response HTTP under load, which provides additional context for choosing gRPC as the baseline communication layer \cite{restgrpccomparison,lahtevenoja2021communication}.

Scalability in containerised streaming systems has been studied through systematic mapping and empirical evaluations. A mapping study on performance analysis techniques for streaming workloads in containerised microservices identifies network contention, scheduling overhead, and coordination costs as primary scalability bottlenecks \cite{systematicstreaming}. Empirical studies of streaming pipelines deployed on Kubernetes further show that streaming approaches scale more gracefully than request-response models when concurrency increases, provided that backpressure and flow control are properly configured \cite{starck2023streaminglb}.

However, most existing evaluations assume static microservice topologies and fixed execution paths. Prior performance studies rarely consider adaptive middleware that dynamically constructs execution chains based on policy evaluation. As a result, there is limited empirical insight into how streaming communication scales when embedded within dynamically adaptive, policy-driven systems such as DYNAMOS. This lack of combined evaluation of streaming, scalability, and adaptivity motivates the experimental focus of this thesis.

\section{Research Gap}\label{sec:research-gap}

A substantial body of work addresses streaming technologies in distributed systems. Many studies focus on performance characteristics of stream processing frameworks such as Apache Spark, Flink, and Storm~\cite{karimov2018benchmarking,fragkoulis2024survey,lote2022realtime}. However, most research on comparing different communication approaches evaluates either traditional request-response patterns or compares streaming frameworks in isolation. When work explicitly compares brokers with RPC mechanisms, it often focuses on unary request-response communication rather than streaming semantics~\cite{ibrahim2026brokers,nuikka2021cloudnative,johansson2020rpc5g,coviello2022dataxc}. No existing work found for this thesis provides a direct empirical comparison between gRPC streaming and broker-based streaming that accounts for long-lived connections, flow control, and backpressure handling in realistic microservice deployments.

Architectural integration and policy enforcement in streaming contexts are also addressed by relatively little prior work. DYNAMOS represents a policy-aware, adaptive approach to microservice composition. Stutterheim~\cite{stutterheim2023thesis} and Stutterheim et al.~\cite{stutterheim2023icsa} describe DYNAMOS and its policy-driven architecture, while related work touches on self-adaptation, policy-aware data exchange, usage control, and gRPC security~\cite{figueira2024selfadaptive,neumaier2021policyaware,park2004ucon,thiyagarajan2025grpcsecurity}. However, these studies do not examine how streaming communication affects policy enforcement, service composition, or middleware behaviour in an adaptive system. The question of how to integrate streaming semantics into policy-driven middleware while preserving its adaptive capabilities therefore remains largely unexplored.

Finally, performance evaluation of streaming systems is well represented in the literature~\cite{karimov2018benchmarking,vogel2024faultRecovery,systematicstreaming,sodankoor2026grpcbench,lote2022realtime}, but these studies tend to focus on throughput and latency in static stream-processing systems. What is missing is an evaluation of how streaming communication affects microservice architectures where services are dynamically composed into policy-approved execution paths. The interaction between streaming semantics, service replication, load balancing, and adaptive middleware components motivates the empirical focus of this thesis.
